% Public text-only snapshot of the instructor's Module 6 autoencoder slides.
% Upstream: /Users/setup/Library/CloudStorage/Box-Box/Teaching/6050/LaTeX/Module 6-AE/slides.tex
% Upstream SHA-256: 46dd07cb1d1d6dae02811865be09af4f28ee37ccd9874d30e38c729ed3a2f75a
% Local preamble dependencies, TikZ source, executable code, external assets,
% private links, assignment links, product-current claims, and claims without a
% sufficiently clear reuse or verification boundary are omitted. The retained
% text and equations document provenance; the book independently derives its
% mathematics, implementations, experiments, and figures.
\documentclass{article}
\usepackage{amsmath}
\usepackage{amssymb}
\usepackage[margin=1in]{geometry}
\usepackage{hyperref}

\title{Autoencoders: Learning Explicit Representations}
\author{Heman Shakeri}
\date{}

\begin{document}
\maketitle

\section*{Public Snapshot Boundary}

This text-only snapshot preserves the instructor-authored lecture spine and selected
equations. Diagram source, executable examples, external assets, private links, and
assignment-specific material have been removed. Remaining technical claims are
historical lecture notes and must be checked independently before book publication.

\section{Motivation: From Classification to Representation}

MLPs and CNNs learn internal features in service of a label. Their representations are
a byproduct of classification rather than the stated learning target. The new question
is whether a network can learn the structure of the data itself without labels.

Explicit representations support compression, anomaly detection, generation, and
transfer. The encoder--decoder pattern also provides a bridge to variable-length
sequence models and translation.

\section{Latent Space and Representation Learning}

Let the observed input be $\mathbf{x}$ and the explicit representation be
$\mathbf{z}$. For compression,

\[
\dim(\mathbf{z}) < \dim(\mathbf{x}).
\]

The motivating hypothesis is that high-dimensional observations may concentrate near
lower-dimensional structure in the ambient space. A representation learner attempts
to capture that structure in coordinates useful to a decoder or downstream task.

% [The original deck's TikZ manifold illustration is omitted.]
% [Unqualified claims that continuous distance automatically equals semantic
% similarity are omitted; the book states the required metric and training assumptions.]

\section{The Autoencoder Architecture}

The autoencoder contract is:

\begin{enumerate}
\item take $\mathbf{x}\in\mathbb{R}^{d}$;
\item encode it as $\mathbf{z}=f_{\theta}(\mathbf{x})\in\mathbb{R}^{k}$;
\item decode it as $\hat{\mathbf{x}}=g_{\phi}(\mathbf{z})\in\mathbb{R}^{d}$.
\end{enumerate}

The training target is the input itself. A narrow latent representation constrains the
information path and asks the network to retain what is useful for reconstruction.

% [The original bottleneck diagram and illustrative latent-width curve are omitted.]
% [The original categorical claim that a bottleneck cannot memorize is omitted. A
% finite training set can still be memorized by a sufficiently expressive encoder and
% decoder; the book treats bottleneck width as a constraint, not a certificate.]
% [The original PyTorch listing is omitted; the book implementation is independent.]

\section{Connection to PCA}

A linear autoencoder with no activation functions has reconstruction

\[
\hat{\mathbf{x}}=W_2W_1\mathbf{x}.
\]

For centered data, an undercomplete linear autoencoder trained with squared
reconstruction loss can learn the same principal subspace as PCA. PCA obtains that
subspace analytically; the autoencoder obtains it through optimization.

For a data matrix $X$, the rank-$k$ reconstruction problem is

\[
\min_{\operatorname{rank}(W_2W_1)\leq k}
\left\|X-W_2W_1X\right\|_F^2.
\]

The optimal reconstruction spans the top-$k$ eigenspace of the covariance matrix.
The factors $W_1$ and $W_2$ are not individually unique or necessarily orthonormal;
latent coordinates can rotate while the reconstructed subspace stays the same.

Adding nonlinearities and depth removes the restriction to one flat subspace. This is
the lecture's bridge from PCA to nonlinear autoencoders and curved data geometry.

% [The PCA and nonlinear-manifold TikZ figures are omitted.]

\section{Reconstruction Objectives}

The assumed conditional observation model determines the reconstruction loss. Under
an isotropic Gaussian decoder,

\[
p(X\mid Z;\phi)=\mathcal{N}\!\left(
X\mid g_{\phi}(Z),\sigma^2I
\right),
\]

maximizing conditional log likelihood is equivalent, up to constants and scale, to
minimizing

\[
\mathcal{L}_{\mathrm{MSE}}(X,\hat X)=\left\|X-\hat X\right\|_2^2.
\]

For Bernoulli observations, the corresponding negative log likelihood is binary cross
entropy:

\[
\mathcal{L}_{\mathrm{BCE}}
=-\sum_i\left[x_i\log \hat x_i+(1-x_i)\log(1-\hat x_i)\right].
\]

% [The original InfoMax derivation and categorical loss-selection prescriptions are
% omitted pending an independent statement of their assumptions.]

\section{Autoencoder Variants}

The same encoder--decoder structure supports different learning contracts:

\begin{itemize}
\item an undercomplete autoencoder constrains latent width;
\item a sparse autoencoder penalizes latent activity;
\item a denoising autoencoder receives corrupted input and reconstructs a clean target;
\item a contractive autoencoder penalizes sensitivity of the representation;
\item a variational autoencoder introduces a stochastic latent variable and a prior.
\end{itemize}

For a denoising autoencoder, let
$\tilde{\mathbf{x}}=\mathbf{x}+\boldsymbol{\epsilon}$ with
$\boldsymbol{\epsilon}\sim\mathcal{N}(0,\sigma^2I)$. A basic objective is

\[
\mathcal{L}_{\mathrm{DAE}}
=\left\|\mathbf{x}-g_{\phi}\!\left(f_{\theta}(\tilde{\mathbf{x}})\right)\right\|_2^2.
\]

The target is still constructed from the input. The architecture remains the same;
what changes is the corruption and objective. This ``same architecture, different
training'' progression connects reconstruction, denoising, and later generative models.

\section{Applications and Convolutional Autoencoders}

The lecture names compression, retrieval, anomaly detection, denoising, and downstream
feature learning as uses of autoencoder representations. For images, convolutional
encoders trade spatial resolution for channels, while decoders restore spatial
resolution. Skip-connected encoder--decoder architectures can pass fine spatial detail
around a bottleneck.

% [Convolutional architecture diagrams are omitted.]
% [The external UMAP asset, its private generation link, and the course-assignment
% notebook link are omitted.]
% [All executable and transposed-convolution teaching code is omitted. No D2L-linked
% material from the associated coding lecture is included in this snapshot.]

\end{document}
